\documentclass[tikz,multi=tikzpicture,border=7pt]{standalone}
\usepackage{ctex}
\usepackage{amsmath,amssymb,bm}
\usetikzlibrary{arrows.meta,calc,decorations.pathmorphing,patterns,positioning,angles,quotes,3d}
\tikzset{
  vec/.style={-{Latex[length=2.6mm]},very thick,blue!72!black},
  force/.style={-{Latex[length=2.6mm]},very thick,red!72!black},
  guide/.style={gray!68,dashed},
  body/.style={draw,very thick,fill=blue!8},
  mass/.style={draw,very thick,fill=orange!45,rounded corners=1pt},
  dim/.style={<->,>=Latex,thin},
  current/.style={-{Latex[length=2.2mm]},very thick,red!70!black},
  field/.style={-{Latex[length=2.2mm]},thick,blue!70!black},
  wave/.style={blue!68!black,thick}
}
\tikzset{vec/.style={-{Latex[length=2.4mm]},very thick,blue!70!black},force/.style={-{Latex[length=2.4mm]},very thick,red!70!black},impulse/.style={force},ray/.style={-{Latex[length=2.2mm]},thick,orange!80!black},wave/.style={blue!70!black,thick},body/.style={draw,very thick,fill=blue!7},guide/.style={gray!65,dashed},dim/.style={<->,>=Latex,thin},resistor/.style={decorate,decoration={zigzag,segment length=4pt,amplitude=2pt},thick}}
\begin{document}

% 1 question-01
\begin{tikzpicture}[x=1cm,y=1cm,font=\small]

\begin{scope}[shift={(-3.8,0)}]
\draw[very thick](-2.8,2)--(-1.6,2)--(-.75,.1)--(.75,.1)--(1.6,2)--(2.8,2);
\draw[very thick,fill=white](0,0)ellipse(.95 and .25);\node[left]at(-.8,.1){$D$};\node[right]at(.8,.1){$E$};\node at(0,.75){同一个上环，质量 $m$};
\draw[very thick](-.75,.1)--(0,-1.8)--(.75,.1);\draw[very thick,fill=white](0,-2.05)circle(.25)node[right=.3cm]{$C,m$};
\draw[guide](-1.6,2)--(-1.6,.9);\draw(-1.6,1.5)arc(-90:-66:.5)node[below right]{$\alpha$};
\node[align=center]at(0,-3){下环：$2T_0=mg$\\上环承受自重及两段绳的向下拉力};\end{scope}
\begin{scope}[shift={(3.5,0)}]\draw[body](0,0)circle(1.15);\draw[very thick](0,-1.15)--(0,-2.3)node[below]{$T_0$};\draw[very thick](0,-1.15)arc(-90:25:1.15);\draw[force](1.04,.486)--++(-.5,1.05)node[above]{$T$};\node at(0,-3){绕行角 $\alpha$；$T/T_0\le e^{\mu\alpha}$};\end{scope}
\end{tikzpicture}

% 2 solution-01
\begin{tikzpicture}[font=\small,x=1cm,y=1cm]
  \draw[->] (-.2,0)--(6.7,0) node[right] {$\alpha$};
  \draw[->] (0,-.2)--(0,8.6) node[above] {$y$};
  \draw[blue!72!black,very thick,domain=0:1.48,samples=120] plot ({4*\x},{4*cos(\x r)}) node[right] {$y_1=\cos\alpha$};
  \draw[orange!82!black,very thick,domain=0:1.48,samples=120] plot ({4*\x},{8*exp(-1.5*\x)}) node[below right] {$y_2=2e^{-\mu\alpha}$};
  \draw[guide] ({4*.5820},0)--({4*.5820},{4*cos(.5820 r)});
  \draw[guide] ({4*1.267},0)--({4*1.267},{4*cos(1.267 r)});
  \node[below] at ({4*.5820},0) {$\alpha_{\min}$};
  \node[below] at ({4*1.267},0) {$\alpha_{\max}$};
  \node[align=left,anchor=west] at (7.2,3.75) {$T_{\min}=T_0e^{-\mu\alpha}$\\[2mm]$T_{\max}=T_0e^{\mu\alpha}$};
  \node[align=left,anchor=west] at (7.2,1.8) {临界相切：\\$\tan\alpha_c=\mu_c$\\$e^{\mu_c\alpha_c}\cos\alpha_c=2$};
\node at(5,-1.5){上环平衡：$2T\cos\alpha=mg+2T_0=2mg$};\end{tikzpicture}

% 3 question-02
\begin{tikzpicture}[x=1cm,y=1cm,font=\small]

\draw[very thick](-4,-1.2)--(4,-1.2);\draw[very thick](-4,1.2)--(4,1.2);\draw[body](1,0)circle(1.2);
\fill(1,0)circle(2pt)node[right]{$C$};\fill(1,1.2)circle(2pt)node[above]{$P_+$};\fill(1,-1.2)circle(2pt)node[below]{$P_-$};
\draw[dim](4.6,-1.2)--(4.6,1.2)node[midway,right]{$2r$};\node at(-2,1.7){上板绕竖直轴转动 $\Omega$};\node at(-2,-1.8){下板固定};\draw[->](1,0)--(3.3,0)node[right]{$x$};\draw[->](1,0)--(1,2.5)node[above]{$z$};\node at(0,-2.6){侧视：上下接触点确在球面直径两端};
\end{tikzpicture}

% 4 solution-02
\begin{tikzpicture}[font=\small,x=1cm,y=1cm]
  \begin{scope}[shift={(-3.8,0)}]
    \draw[->] (-2.4,0)--(2.7,0) node[right] {$x$}; \draw[->] (0,-2.4)--(0,2.7) node[above] {$y$};
    \draw[very thick,blue!70!black] (0,0) circle (1.7);
    \fill (1.7,0) circle (2pt) node[above right] {$C(0)=(l,0)$};
    \fill (-1.7,0) circle (2pt) node[left] {$C(2\pi/\Omega)=(-l,0)$};
    \draw[current] (1.2,1.2) arc (45:210:1.7) node[midway,left] {$\Omega/2$};
    \node at (0,-3.0) {球心的半速圆周运动};
  \end{scope}
  \begin{scope}[shift={(4.1,0)}]
    \draw[->] (0,-2.4)--(0,2.7) node[above] {$\bm e_z$};
    \draw[->] (0,0)--(-2.1,0) node[left] {$\bm e_r$};
    \draw[vec] (0,0)--(-2.1,-1.25) node[below left] {$\bm\omega'=-(\Omega/2r)\bm r_C-(\Omega/2)\bm e_z$};
    \draw (0,-.7) arc (-90:-149:.7) node[midway,below] {$\chi$};
    \node[align=center] at (0,-3.15) {$|\bm\omega'|\,\dfrac{2\pi}{\Omega}=2\pi n$\\[1mm]$\displaystyle \frac rl=\frac1{\sqrt{4n^2-1}}$};
  \end{scope}
\end{tikzpicture}

% 5 question-03
\begin{tikzpicture}[x=1cm,y=1cm,font=\small]
\begin{scope}[shift={(-4,0)}]\draw[gray!80,thick] plot coordinates {(-3.00000,0.00000) (-2.50000,-0.86603)};
\draw[gray!80,thick] plot coordinates {(-3.00000,0.00000) (-2.00000,0.00000)};
\draw[gray!80,thick] plot coordinates {(-3.00000,1.73205) (-2.00000,1.73205)};
\draw[gray!80,thick] plot coordinates {(-3.00000,1.73205) (-2.50000,2.59808)};
\draw[gray!80,thick] plot coordinates {(-2.00000,1.73205) (-2.50000,2.59808)};
\draw[gray!80,thick] plot coordinates {(-2.00000,1.73205) (-1.50000,0.86603)};
\draw[gray!80,thick] plot coordinates {(-2.00000,1.73205) (-1.00000,1.73205)};
\draw[gray!80,thick] plot coordinates {(-3.00000,-1.73205) (-2.00000,-1.73205)};
\draw[gray!80,thick] plot coordinates {(-3.00000,-1.73205) (-2.50000,-0.86603)};
\draw[gray!80,thick] plot coordinates {(-2.00000,-1.73205) (-2.50000,-0.86603)};
\draw[gray!80,thick] plot coordinates {(-2.00000,-1.73205) (-1.50000,-2.59808)};
\draw[gray!80,thick] plot coordinates {(-2.00000,-1.73205) (-1.00000,-1.73205)};
\draw[gray!80,thick] plot coordinates {(-2.50000,-0.86603) (-2.00000,0.00000)};
\draw[gray!80,thick] plot coordinates {(-2.00000,0.00000) (-1.00000,0.00000)};
\draw[gray!80,thick] plot coordinates {(-2.00000,0.00000) (-1.50000,0.86603)};
\draw[gray!80,thick] plot coordinates {(-1.00000,0.00000) (-1.50000,0.86603)};
\draw[gray!80,thick] plot coordinates {(-1.00000,0.00000) (-0.50000,-0.86603)};
\draw[gray!80,thick] plot coordinates {(-1.00000,0.00000) (0.00000,0.00000)};
\draw[gray!80,thick] plot coordinates {(-1.50000,0.86603) (-1.00000,1.73205)};
\draw[gray!80,thick] plot coordinates {(-1.00000,1.73205) (0.00000,1.73205)};
\draw[gray!80,thick] plot coordinates {(-1.00000,1.73205) (-0.50000,2.59808)};
\draw[gray!80,thick] plot coordinates {(0.00000,1.73205) (-0.50000,2.59808)};
\draw[gray!80,thick] plot coordinates {(0.00000,1.73205) (0.50000,0.86603)};
\draw[gray!80,thick] plot coordinates {(0.00000,1.73205) (1.00000,1.73205)};
\draw[gray!80,thick] plot coordinates {(-1.50000,-2.59808) (-1.00000,-1.73205)};
\draw[gray!80,thick] plot coordinates {(-1.00000,-1.73205) (0.00000,-1.73205)};
\draw[gray!80,thick] plot coordinates {(-1.00000,-1.73205) (-0.50000,-0.86603)};
\draw[gray!80,thick] plot coordinates {(0.00000,-1.73205) (-0.50000,-0.86603)};
\draw[gray!80,thick] plot coordinates {(0.00000,-1.73205) (0.50000,-2.59808)};
\draw[gray!80,thick] plot coordinates {(0.00000,-1.73205) (1.00000,-1.73205)};
\draw[gray!80,thick] plot coordinates {(-0.50000,-0.86603) (0.00000,0.00000)};
\draw[gray!80,thick] plot coordinates {(0.00000,0.00000) (1.00000,0.00000)};
\draw[gray!80,thick] plot coordinates {(0.00000,0.00000) (0.50000,0.86603)};
\draw[gray!80,thick] plot coordinates {(1.00000,0.00000) (0.50000,0.86603)};
\draw[gray!80,thick] plot coordinates {(1.00000,0.00000) (1.50000,-0.86603)};
\draw[gray!80,thick] plot coordinates {(1.00000,0.00000) (2.00000,0.00000)};
\draw[gray!80,thick] plot coordinates {(0.50000,0.86603) (1.00000,1.73205)};
\draw[gray!80,thick] plot coordinates {(1.00000,1.73205) (2.00000,1.73205)};
\draw[gray!80,thick] plot coordinates {(1.00000,1.73205) (1.50000,2.59808)};
\draw[gray!80,thick] plot coordinates {(2.00000,1.73205) (1.50000,2.59808)};
\draw[gray!80,thick] plot coordinates {(2.00000,1.73205) (2.50000,0.86603)};
\draw[gray!80,thick] plot coordinates {(2.00000,1.73205) (3.00000,1.73205)};
\draw[gray!80,thick] plot coordinates {(0.50000,-2.59808) (1.00000,-1.73205)};
\draw[gray!80,thick] plot coordinates {(1.00000,-1.73205) (2.00000,-1.73205)};
\draw[gray!80,thick] plot coordinates {(1.00000,-1.73205) (1.50000,-0.86603)};
\draw[gray!80,thick] plot coordinates {(2.00000,-1.73205) (1.50000,-0.86603)};
\draw[gray!80,thick] plot coordinates {(2.00000,-1.73205) (2.50000,-2.59808)};
\draw[gray!80,thick] plot coordinates {(2.00000,-1.73205) (3.00000,-1.73205)};
\draw[gray!80,thick] plot coordinates {(1.50000,-0.86603) (2.00000,0.00000)};
\draw[gray!80,thick] plot coordinates {(2.00000,0.00000) (3.00000,0.00000)};
\draw[gray!80,thick] plot coordinates {(2.00000,0.00000) (2.50000,0.86603)};
\draw[gray!80,thick] plot coordinates {(3.00000,0.00000) (2.50000,0.86603)};
\draw[gray!80,thick] plot coordinates {(2.50000,0.86603) (3.00000,1.73205)};
\draw[gray!80,thick] plot coordinates {(2.50000,-2.59808) (3.00000,-1.73205)};
\fill (-3.0000,0.0000) circle(1.4pt);
\fill (-3.0000,1.7321) circle(1.4pt);
\fill (-2.0000,1.7321) circle(1.4pt);
\fill (-2.5000,2.5981) circle(1.4pt);
\fill (-3.0000,-1.7321) circle(1.4pt);
\fill (-2.0000,-1.7321) circle(1.4pt);
\fill (-2.5000,-0.8660) circle(1.4pt);
\fill (-2.0000,0.0000) circle(1.4pt);
\fill (-1.0000,0.0000) circle(1.4pt);
\fill (-1.5000,0.8660) circle(1.4pt);
\fill (-1.0000,1.7321) circle(1.4pt);
\fill (0.0000,1.7321) circle(1.4pt);
\fill (-0.5000,2.5981) circle(1.4pt);
\fill (-1.5000,-2.5981) circle(1.4pt);
\fill (-1.0000,-1.7321) circle(1.4pt);
\fill (0.0000,-1.7321) circle(1.4pt);
\fill (-0.5000,-0.8660) circle(1.4pt);
\fill (0.0000,0.0000) circle(1.4pt);
\fill (1.0000,0.0000) circle(1.4pt);
\fill (0.5000,0.8660) circle(1.4pt);
\fill (1.0000,1.7321) circle(1.4pt);
\fill (2.0000,1.7321) circle(1.4pt);
\fill (1.5000,2.5981) circle(1.4pt);
\fill (0.5000,-2.5981) circle(1.4pt);
\fill (1.0000,-1.7321) circle(1.4pt);
\fill (2.0000,-1.7321) circle(1.4pt);
\fill (1.5000,-0.8660) circle(1.4pt);
\fill (2.0000,0.0000) circle(1.4pt);
\fill (3.0000,0.0000) circle(1.4pt);
\fill (2.5000,0.8660) circle(1.4pt);
\fill (3.0000,1.7321) circle(1.4pt);
\fill (2.5000,-2.5981) circle(1.4pt);
\fill (3.0000,-1.7321) circle(1.4pt);
\foreach\p/\lab in{(0,0)/A,(1,0)/B,(.5,.866)/C}{\fill[red!70!black]\p circle(2.5pt)node[above]{$\lab$};}
\node at(0,-3.2){角点相接的三角网络：每节点四配位};\end{scope}
\begin{scope}[shift={(4,0)}]\coordinate(O)at(0,0);\draw[very thick](-2,1.5)--(O)--(2,1.5);\draw[very thick](O)--(0,-1.7);\node[above]at(-2,1.5){$A$};\node[above]at(2,1.5){$B$};\node[below]at(0,-1.7){$C$};\node[above left]at(-1,.8){$Z_A$};\node[above right]at(1,.8){$Z_B$};\node[right]at(0,-.8){$Z_C$};\node at(0,-3.2){完整网络：$Z_A=Z_B=Z_C=r/4$};\end{scope}
\end{tikzpicture}

% 6 solution-03
\begin{tikzpicture}[x=1cm,y=1cm,font=\small]

\begin{scope}[shift={(-3.6,0)}]\coordinate(O)at(0,0);\draw[very thick](-2,1.5)--node[below left]{$r/2$}(O)--node[below right]{$r/2$}(2,1.5);\draw[very thick](O)--(0,-1.6)node[midway,right]{$r/8$};
\node[above]at(-2,1.5){$A$};\node[above]at(2,1.5){$B$};\node[below]at(0,-1.6){$C$};\draw[guide](-2,1.5)--(2,1.5);\draw[red!70!black,very thick](-.2,1.3)--(.2,1.7);\node at(0,2.2){$AB$ 直连电阻已断开};\node at(0,-2.5){$R'_{AB}=r$，$R'_{AC}=R'_{BC}=5r/8$};\end{scope}
\begin{scope}[shift={(3.7,0)}]\draw[->](-2.5,0)--(2.7,0)node[right]{$\omega t$};
\foreach\p/\sty in{0/solid,120/dashed,240/dotted}{\draw[wave,\sty,domain=-2.4:2.4,samples=80]plot(\x,{1.4*cos(90*\x+\p)});}
\node[align=center]at(0,-2.5){完整网络的三相驱动\\$u_A+u_B+u_C=0$\\$I_A=(4u_0/r)\cos\omega t$};\end{scope}
\end{tikzpicture}

% 7 question-04
\begin{tikzpicture}[font=\small,x=1cm,y=1cm]
  \begin{scope}[shift={(-3.8,0)}]
    \draw[very thick] (0,0) circle (2.2); \fill (-2.2,0) circle (3pt) node[left] {$A$}; \fill (0,0) circle (2pt) node[above] {$O$};
    \coordinate (P) at (1.25,-1.81); \shade[ball color=orange!55] (P) circle (.24) node[below right] {$m,q$};
    \draw[very thick] (-2.2,0)--(P) node[midway,below] {$r$};
    \draw[guide] (0,0)--(P); \draw[current] (1.65,1.45) arc (40:110:.8) node[right] {$\omega$};
    \node at (.3,.85) {$\otimes\,\bm B$};
    \node at (0,-3.0) {带电小珠在铰接绝缘环上滑动};
  \end{scope}
  \begin{scope}[shift={(4.0,0)}]
    \draw[very thick] (0,0) circle (2.2); \fill (-2.2,0) circle (3pt) node[left] {$A$}; \fill (0,0) circle (2pt) node[right] {$O$};
    \draw[very thick] (-1.1,0) circle (1.1); \fill (-1.1,0) circle (2pt) node[above] {$O'$};
    \shade[ball color=orange!55] (-2.2,0) circle (.22) node[below] {$q$};
    \draw[current] (1.7,1.35) arc (35:105:.8) node[right] {$\omega_0'$};
    \node at (.45,.8) {$\otimes\,\bm B$};
    \node at (0,-3.0) {小环在大环内纯滚动};
  \end{scope}
\end{tikzpicture}

% 8 solution-04
\begin{tikzpicture}[x=1cm,y=1cm,font=\small]

\begin{scope}[shift={(-3.6,0)}]\draw[body](0,0)circle(2);\fill(-2,0)circle(2pt)node[left]{$A$};\fill(.9,-1.786)circle(3pt)node[below right]{$q,m$};\draw[very thick](-2,0)--(.9,-1.786)node[midway,below]{$r$};\node at(0,.4){$\otimes\,B$};\node at(0,1.2){$\odot\,L_z$};\node[align=center]at(0,-3){$L_z-\tfrac12qBr^2=$ 常量\\磁场不做功};\end{scope}
\begin{scope}[shift={(3.4,0)}]\coordinate(A)at(0,0);\coordinate(O)at(2.2,0);\draw[very thick](O)circle(2.2);\coordinate(Op)at(1.56907,0.90107);\draw[very thick](Op)circle(1.1);\fill(0.93813,0)circle(3pt)node[below]{$q$};
\fill(A)circle(2pt)node[left]{$A$};\fill(O)circle(2pt)node[right]{$O$};\fill(Op)circle(2pt)node[above]{$O'$};\draw[guide](O)--(A);\draw[very thick](O)--(Op);
\draw(1.55,0)arc(180:125:.65)node[midway,left]{$\alpha$};\node[align=center]at(2,-3){图取 $\alpha=55^\circ$；小环确在大环内\\小珠在 $AO$ 上：$r=R(1-\cos\alpha)$\\$\beta=2\alpha$ 是相对接触半径的夹角};\end{scope}
\end{tikzpicture}

% 9 question-05
\begin{tikzpicture}[x=1cm,y=1cm,font=\small]

\fill(-4,0)circle(3pt)node[below]{$Q$};\coordinate(C)at(1,0);\draw[vec](-4,0)--(C)node[midway,above]{$\bm r$（从 $Q$ 指向质心）};
\draw[very thick](-.0,-1.0)--(2,1);\foreach\p/\lab in{(0,-1)/{-q},(1,0)/{2q},(2,1)/{-q}}{\fill\p circle(3pt)node[right]{$\lab$};}
\draw[vec](C)--(2.3,1.3)node[above]{$\hat{\bm n}$};\draw[guide](C)--(3.1,0);\draw(1.8,0)arc(0:45:.8)node[right]{$\theta$};\draw[dim](.2,-1.2)--(1.2,-.2)node[midway,right]{$l$};\draw[dim](1.2,-.2)--(2.2,.8)node[midway,right]{$l$};\node at(0,-2){线性三电荷分子，总电荷与偶极矩均为零};
\end{tikzpicture}

% 10 solution-05
\begin{tikzpicture}[x=1cm,y=1cm,font=\small]

\begin{scope}[shift={(-3.5,0)}]\draw[->](-2.2,0)--(2.3,0)node[right]{$\hat{\bm r}$};\draw[vec](0,0)--(1.3,1.3)node[above]{$\hat{\bm n}$};\draw[->](0,0)--(0,-1.7)node[below]{$\hat{\bm\theta}$};\draw(1,0)arc(0:45:1)node[right]{$\theta$};
\node[align=center]at(0,-3){$U=K(1-3\cos^2\theta)/r^3$\\$F_r=3K(1-3\cos^2\theta)/r^4$\\$F_\theta=-6K\sin\theta\cos\theta/r^4$};\end{scope}
\begin{scope}[shift={(3.7,0)}]\draw[very thick](-1.5,0)--(1.5,0);\foreach\x/\lab in{-1.5/{-q},0/{2q},1.5/{-q}}{\fill(\x,0)circle(3pt)node[below]{$\lab$};}\node at(0,1.6){$\odot\,\bm b$（垂直分子）};
\draw[force](-1.5,0)--(-1.5,-1);\draw[force](1.5,0)--(1.5,1);\node[align=center]at(0,-3){两端力相反、力矩同向\\$\bm\tau=ql^2[\bm b-\hat{\bm n}(\hat{\bm n}\cdot\bm b)]$\\$\bm b\parallel\hat{\bm n}$ 时各端电场为零};\end{scope}
\end{tikzpicture}

% 11 question-06
\begin{tikzpicture}[x=1cm,y=1cm,font=\small]

\fill[blue!8](-4,0)rectangle(3.7,3);\draw[very thick](-4,0)--(3.7,0)node[right]{岸线};\draw[->](-4,0)--(4.2,0)node[right]{$X=x/y_0$};\draw[->](0,0)--(0,3.5)node[above]{$Y=y/y_0$};\draw[guide](-4,2.5)--(3.7,2.5);\fill(0,2.5)circle(2pt)node[above left]{$A$};\draw[vec](.2,2.5)--(1.9,2.5)node[right]{$v_b=2v_0$};
\draw[wave,domain=0:.523599,samples=70]plot({2.5*\x},{5*cos((\x+pi/3)r)});\fill({2.5*pi/6},0)circle(2pt)node[below]{$B$};\node[align=left]at(-2,1.3){$h(y)=4h_0/(5-Y^2)$\\$v_0=\sqrt{gh_0}$};\node at(0,-.9){射线仅画 $A$ 到岸线 $B$ 的有效段};
\end{tikzpicture}

% 12 solution-06
\begin{tikzpicture}[x=1cm,y=1cm,font=\small]

\begin{scope}[shift={(-3.7,0)}]\draw[->](-.3,0)--(3,0)node[right]{$X$};\draw[->](0,-.2)--(0,3.2)node[above]{$Y$};\draw[wave,domain=0:pi/6,samples=80]plot({4*\x},{6*cos((\x+pi/3)r)});\fill(0,3)circle(2pt)node[left]{$A:(0,1)$};\fill({2*pi/3},0)circle(2pt)node[below]{$B:(\pi/6,0)$};\node at(1.2,-1){射线 $Y=2\cos(X+\pi/3)$};\end{scope}
\begin{scope}[shift={(3.7,0)}]\draw[->](-3.5,0)--(1,0)node[right]{$X$};\draw[->](0,-.2)--(0,3.2)node[above]{$Y$};\draw[orange!80!black,very thick,domain=pi/3:pi/2,samples=100,variable=\p]plot({2.6*(-2*(\p-pi/3)+sin(2*\p r)-sqrt(3)/2)},{6*cos(\p r)});\fill({2.6*(-.9237395)},1.5)circle(2pt)node[left]{$(-.92374,.5)$};\draw[guide](-3,1.5)--(.8,1.5);\node at(-1,-1){同一时刻的波前，不是单条射线};\end{scope}
\end{tikzpicture}

% 13 question-07
\begin{tikzpicture}[font=\small,x=1cm,y=1cm]
  \draw[very thick] (-2.3,-3.5)--(-2.3,2.5)--(1.1,2.5)--(1.1,-.7);
  \draw[very thick] (-2.3,-3.5)--(-1.65,-3.5);\draw[very thick] (-1.65,-3.5)--(-1.65,1.85)--(.45,1.85)--(.45,-.7);
  \draw[fill=gray!55,draw=none] (-2.29,1.55) rectangle (-1.66,2.48);
  \node[left] at (-2.3,2.5) {$A$}; \node[right] at (1.1,2.5) {$B$};
  \draw[dim] (-3.0,-3.5)--(-3.0,2.5) node[midway,left] {$l_1$};
  \draw[dim] (-2.3,3.05)--(1.1,3.05) node[midway,above] {$l_2$};
  \draw[dim] (1.8,-.7)--(1.8,2.5) node[midway,right] {$l_3$};
  \draw[dim] (-1.25,1.55)--(-1.25,2.48) node[midway,right] {$l_4$};
  \node[align=left,anchor=west] at (3.0,1.5) {$l_1=100\,\mathrm{cm}$\\$l_2=5\,\mathrm{cm}$\\$l_3=50\,\mathrm{cm}$\\$l_4=10\,\mathrm{cm}$};
  \node[align=left,anchor=west] at (3.0,-1.0) {左端封闭\\右端通大气\\灰色为水银柱};
\end{tikzpicture}

% 14 solution-07
\begin{tikzpicture}[x=1cm,y=1cm,font=\small]

\begin{scope}[shift={(-3.7,0)}]\draw[very thick](0,-1.8)--(0,1.3)--(1.5,1.3)--(1.5,-.7);\draw[very thick](.4,-1.8)--(.4,.9)--(1.1,.9)--(1.1,-.7);\draw[very thick](0,-1.8)--(.4,-1.8);
\fill[gray!60](.01,.35)rectangle(.39,1.29);\node[left]at(0,1.3){$A$};\node[right]at(1.5,1.3){$B$};\node at(.7,-2.5){初态 $q_0=90$ cm，$p_i=86$ cmHg};\end{scope}
\begin{scope}[shift={(3.7,0)}]\draw[->](-.4,0)--(5.9,0)node[right]{$q$/cm};\draw[->](0,-.4)--(0,3.5)node[above]{$p$/cmHg};
\draw[wave,domain=90:155,samples=100]plot({(\x-90)/12},{(7740/\x-50)/12});\draw[orange!80!black,very thick] (0,3)--(.417,2.583)--(.833,1.75)--(1.25,1.333)--(4.583,1.333)--(5.417,2.167);
\node[below]at(0,0){90};\node[below]at(1.25,0){105};\node[below]at(4.583,0){145};\node[align=center]at(2.7,-1.4){橙：实际分段平衡压力 $p_{\rm eq}$\\蓝：$p_g=7740/q$，$T=300$ K\\正向初扰动失稳；弱阻尼后停在稳定交点};\end{scope}
\end{tikzpicture}

% 15 question-08
\begin{tikzpicture}[x=1cm,y=1cm,font=\small]

\draw[->](-4.5,0)--(4.7,0)node[right]{$x$};\draw[body](-3,0)circle(.25)node[below=.25cm]{电子 $\gamma m$};\draw[vec](-2.5,0)--(-1,0)node[above]{$v$};\draw[body](0,0)circle(.25)node[below=.25cm]{静止正电子 $m$};
\draw[wave,very thick](0,0)--(2.3,1.4)node[right]{$h\nu_1$};\draw[wave,very thick](0,0)--(2.3,-1.4)node[right]{$h\nu_2$};\draw(.8,-.49)arc(-31:31:.94)node[right]{$\theta$};\node at(0,-2.2){出射光子总动量沿入射电子方向};
\end{tikzpicture}

% 16 solution-08
\begin{tikzpicture}[x=1cm,y=1cm,font=\small]

\begin{scope}[shift={(-3.5,0)}]\draw[->](-.5,0)--(3.3,0)node[right]{$p_x$};\draw[->](0,-2.1)--(0,2.1)node[above]{$p_y$};
\draw[vec](0,0)--(1.5,1.5)node[above]{$\bm p_1$};\draw[vec](0,0)--(1.5,-1.5)node[below]{$\bm p_2$};\draw[vec,red!70!black](0,0)--(3,0)node[above]{$\bm p_{\rm tot}$};\draw[guide](1.5,1.5)--(3,0)--(1.5,-1.5);
\node[align=center]at(1.5,-3){示例 $\gamma=3$，等能光子夹角 $90^\circ$\\$\bm p_1+\bm p_2=\bm p_{\rm tot}$};\end{scope}
\begin{scope}[shift={(3.7,0)}]\draw[wave,very thick](0,0)--(0,1.6)node[above]{$E_\gamma'$};\draw[wave,very thick](0,0)--(0,-1.6)node[below]{$E_\gamma'$};\draw[vec](-1.7,-2.6)--(.7,-2.6)node[right]{$\beta_c c$};\node[align=center]at(0,2.3){质心系两光子等能、反向};\node at(0,-3.4){$\theta_{\min}=2\arccos\sqrt{(\gamma-1)/(\gamma+1)}$};\end{scope}
\end{tikzpicture}

\end{document}
